\documentclass[12pt, a4paper]{article}

% Packages for formatting and functionality
\usepackage[utf8]{inputenc}
\usepackage[T1]{fontenc}
\usepackage{geometry}
\geometry{a4paper, margin=1in}
\usepackage{booktabs} % For professional tables
\usepackage{tabularx}
\usepackage{hyperref} % For clickable links and TOC
\usepackage{xcolor}
\usepackage{graphicx}
\usepackage{amsmath}
\usepackage{enumitem}
\usepackage{listings}
\usepackage{float}
\usepackage{lmodern}

% Setup Python Code Snippet Styling
\lstset{
    language=Python,
    basicstyle=\ttfamily\footnotesize,
    keywordstyle=\color{blue}\bfseries,
    stringstyle=\color{red},
    commentstyle=\color{green!60!black},
    numbers=left,
    numberstyle=\tiny\color{gray},
    stepnumber=1,
    frame=single,
    breaklines=true,
    backgroundcolor=\color{black!5},
    showstringspaces=false
}

% Setup hyperlink styling
\hypersetup{
    colorlinks=true,
    linkcolor=blue,
    filecolor=magenta,      
    urlcolor=cyan,
    pdftitle={Sequential ATPG Project Report},
}

% Title Configuration
\title{\textbf{SEQUENTIAL ATPG} \\ \Large Extended D-Algorithm vs. 9-Valued Logic (Muth's Algorithm) \\ \vspace{0.5cm} \large A Comprehensive Study and Implementation Guide for Sequential Digital Circuits \\ \vspace{0.5cm} \normalsize Group Project \& Viva Preparation Report}
\author{VLSI Testing \& Verification}
\date{Benchmarks: ISCAS-89 s27 and s298}

\begin{document}

\maketitle

\begin{abstract}
\noindent This report presents a comprehensive study and implementation of two Automatic Test Pattern Generation (ATPG) algorithms for sequential digital circuits. The first algorithm is the Extended D-Algorithm, which employs general 5-valued logic (0, 1, X, D, D\_BAR) and is widely used for combinatorial mapping. The second algorithm is Muth's 9-Valued Logic Algorithm, which extends representational power to 9 distinct states, decoupling good and faulty circuit signals to handle unknown flip-flop initialization states. 

The project evaluates both systems implemented from scratch in Python, utilizing time-frame expansion on the ISCAS-89 benchmark circuitry (\texttt{s27} with 3 flip-flops and \texttt{s298} with 14 flip-flops). Extensive experiments reveal a striking contrast between the algorithms: the 5-valued Extended D-Algorithm generates optimistic but highly \textit{fragile} test patterns. It claims 100\% theoretical fault detection on \texttt{s298} with 8 time frames, but strict fault simulation verification exposes that these tests only yield 95.59\% concrete coverage because of unresolved `X` states. Conversely, Muth's 9-valued algorithm accurately scopes partial states inherently producing highly robust vectors. While it conservatively caps detectability at 89.71\% under identical limits, those test vectors yield an explosive 98.16\% true verified coverage. This mathematical soundness comes strictly at the expense of computational bounds resulting in 3-10x longer path evaluation times.
\end{abstract}

\newpage
\tableofcontents
\newpage

\section{Introduction \& Problem Statement}

\subsection{The Challenge of Testing Digital Circuits}
Modern VLSI (Very Large Scale Integration) chips contain millions of gates. After fabrication, imperfections cause wires to permanently drive 1 or permanently drive 0. The goal of ATPG (Automatic Test Pattern Generation) is determining the exact sequence of primary input assignments required to expose such faults at primary output terminals.

While combinational circuits (stateless) are well resolved via Roth's D-Algorithm, \textbf{sequential circuits} pose a devastating challenge: structural faults depend on state machines trapped inside memory elements (D-type Flip-Flops). Therefore, evaluating a problem intrinsically spans multiple clock cycles where we must first explicitly initialize the flip-flop values.

\subsection{Project Problem Statement}
This architecture tackles these obstacles precisely:
\begin{itemize}
    \item Parsing complex Verilog gate-level netlists with sequential loops mapping to native graph structures.
    \item Systemically testing single Stuck-at-0 (SA0) and Stuck-at-1 (SA1) faults.
    \item Constructing ATPG trees across two independent multi-value engines (5-value vs 9-value).
    \item Implementing robust Time-Frame Expansion (Circuit Unrolling) dictated inherently by Sequential Depth calculations ($d_{\text{seq}}$).
    \item Applying simulated verifications evaluating vector true-coverage dynamically vs raw tree node outputs.
\end{itemize}

\section{Background: Logic Frameworks}

\subsection{D-Calculus and 5-Valued Logic (Mode A)}
The 5-valued system operates under strict combinatorial assumptions mapping both theoretical networks tightly:
\begin{table}[H]
\centering
\caption{5-Valued Logic System}
\begin{tabular}{cccl}
\toprule
\textbf{Value} & \textbf{Good Circuit} & \textbf{Faulty Circuit} & \textbf{Meaning} \\
\midrule
0 & 0 & 0 & Definite 0 in both \\
1 & 1 & 1 & Definite 1 in both \\
X & Unknown & Unknown & Undetermined initialization \\
D & 1 & 0 & Fault effect: good=1, faulty=0 \\
D\_BAR (D')& 0 & 1 & Fault effect: good=0, faulty=1 \\
\bottomrule
\end{tabular}
\end{table}

\subsection{9-Valued Logic: Muth's Algorithm (Mode B)}
Muth (1976) introduced decoupled representations explicitly tracking state behaviors uncoupled from rigid combinations. Values are tracked as \texttt{g/f} sequences. The four revolutionary unknown representations (0/X, X/0, 1/X, X/1) allow Muth's objective solver to isolate variables that strictly matter without risking the failure boundaries that collapse the 5-val ExtD matrix into unresolved loops.

\section{Codebase Architecture: File-by-File Technical Guide}

To fully grasp the implementation, it is vital to understand exactly what each Python file accomplishes sequentially. This section cross-references the core theories directly to the repository line mappings.

\subsection{\texttt{parser.py} (Netlist Interpretation)}
This file is the bedrock. It translates Verilog text files into in-memory object structures suitable for algorithms.
\begin{itemize}
    \item \textbf{Data Structures (Line 20 \& 32):} 
        \begin{itemize}
            \item \texttt{Gate}: Lightweight dataclass defining a component (\texttt{name}, \texttt{gate\_type}, list of \texttt{inputs}, and single \texttt{output}). 
            \item \texttt{Circuit}: The primary network container. It maintains dictionaries of all \texttt{gates}, sets of \texttt{primary\_inputs}/\texttt{primary\_outputs}, and vitally establishes \texttt{fanin} and \texttt{fanout} dependency maps to allow algorithmic traversal in $\mathcal{O}(1)$ time.
        \end{itemize}
    \item \textbf{Main Parser Logic (\texttt{parse\_verilog}, Line 86)}: Uses regex routines scanning standard Verilog (e.g. \texttt{dff}, \texttt{nand}) splitting internal wires.
    \item \textbf{Topological Sorting (\texttt{topological\_order}, Line 174)}: Employs Kahn's BFS algorithm spanning \texttt{in\_degree} tracking mapping gates layer by layer. This sorting bounds that during evaluations, no dependent nodes compute before parent vectors. 
\end{itemize}

\subsection{\texttt{logic5val.py} \& \texttt{logic9val.py} (Gate Combinations)}
These two files determine how values flow across primitives like AND and OR given specific mathematical environments.
\begin{itemize}
    \item \textbf{\texttt{logic5val.py}}: Hardcodes static dictionaries resolving arrays statically (e.g. \texttt{AND(D, D\_BAR) = 0}).
    \item \textbf{\texttt{logic9val.py} Evaluators (Line 64)}: Cleanly splits `good/fault` variables. The loop inside \texttt{and\_9val(*args)} parses inputs via \texttt{decode()} (Line 27), applying basic \texttt{\_and3} values (like $\{0,1,X\}$) strictly linearly before calling \texttt{encode(g, f)} (Line 34).
    \item \textbf{Merge Logic (\texttt{merge\_9val}, Line 170):} Tracks independent bounds iteratively overlapping partial paths. If two paths trace matching data limits like \texttt{0/X} returning to \texttt{X/0}, the intersection strictly binds back returning \texttt{0/0}. If bounds clash fatally, returning \texttt{None} triggering DFS skips.
\end{itemize}

\subsection{\texttt{sgraph.py} (Topological Memory Scaling Limits)}
Calculates dependencies bridging internal DFF memory logic.
\begin{itemize}
    \item \textbf{\texttt{build\_sgraph} (Line 25):} Recursively traces BFS parameters sweeping backwards iteratively from an arbitrary D-pin bounding up combinatorial nets until it strikes a different Q-pin source generating DAG adjacencies organically.
    \item \textbf{\texttt{compute\_sequential\_depth} (Line 97):} The critical function evaluating geometric dependencies natively. Maps level-1 components resolving loops hierarchically extracting $d_{\text{seq}}$, capping algorithmic search loop arrays.
\end{itemize}

\subsection{\texttt{timeframe.py} (Circuit Unroller)}
Translates sequential logic into a purely combinational dimension by copying instances structurally.
\begin{itemize}
    \item \textbf{Unroll Constructor (\texttt{unroll}, Line 72):} Produces an \texttt{UnrolledCircuit} dataclass loop appending \texttt{@t} time bounds specifically onto variable strings resolving cross-contamination natively. 
    \item \textbf{DFF\_CONNECT Links (Line 139):} By traversing mapped frames `t` sequentially backward iteratively, flip keys map DFF bounds linking structural output wires iteratively into immediate frame dependencies strictly without physical memory blocks.
    \item \textbf{Pseudo-Primary Inputs (Line 134):} Extrapolates Q-Outputs mapped inside the earliest TimeFrame loop bounds (-k) as direct external \texttt{PPIs} strictly acting equivalently to configurable PI entry parameters organically.
\end{itemize}

\subsection{\texttt{atpg\_9val.py} (The Muth Search Kernel)}
This file represents the flagship solver, natively relying on recursive bounding to compute variables organically. (Note: \texttt{atpg\_ext\_d.py} evaluates identically structurally aside from value loops mappings).
\begin{itemize}
    \item \textbf{Search Recursion (\texttt{\_search}, Line 144)}: Binds logic recursively trying assignment configurations against structural trees, triggering trace increments incrementally checking limits (`backtrack_limit`).
    \item \textbf{Implication (\texttt{\_imply}, Line 316)}: Resolves DFF logic internally explicitly across bounds continuously traversing layers evaluating `new_vals` updating circuit vectors using independent 9-Val topological sorts. 
    \item \textbf{D-Frontier Extractor (\texttt{\_d\_frontier}, Line 402)}: Sweeps structural DAGs extracting bounded parameters (including frame limits jumping flip boundaries on Line 417). Identifies active target ranges needing specific propagation arrays implicitly routing forward.
    \item \textbf{Backtracing Routines (\texttt{\_backtrace}, Line 239)}: Using objective maps defined strictly via controlling variables boundaries, paths span backward dynamically mapping non-controlling PI variables to fulfill target matrices mapping actively.
\end{itemize}

\subsection{\texttt{fault\_sim.py} \& \texttt{visualize.py}}
\begin{itemize}
    \item \textbf{\texttt{fault\_sim.py} (\texttt{FaultSimulator.simulate}, Line 22):} The core concrete binary verification simulator. Given an array list of target inputs, strictly assigns inputs mapping exclusively 0/1 loops bypassing symbolic arrays `X` bounds entirely confirming whether outputs deviate concretely inside target bounds natively.
    \item \textbf{\texttt{visualize.py} (\texttt{TraceRecorder}, Line 139):} Uses `pyvis.network.Network` rendering shapes mapped across DOM scopes identically. Captures inner frames from \texttt{\_search} scopes allowing HTML navigation routines inside the resulting DOM traces dynamically.
\end{itemize}

\section{Experimental Methodology \& DSE Results}

We evaluated a sweeping configuration test array composed of 16 runs measuring 5-Val ExtD versus 9-Val Muth across s27 and s298. Over the design exploration mapped in `results/design_space.csv`, the structural limits systematically revealed the behaviors.

\subsection{Full Results Array (Extracted from CSV)}

\begin{table}[H]
\centering
\caption{Key Comparative Array (frames=8 limits=10/50)}
\resizebox{\textwidth}{!}{
\begin{tabular}{lcccccc}
\toprule
\textbf{Algorithm} & \textbf{Frames} & \textbf{ATPG Scope Cov} & \textbf{True Verif Cov} & \textbf{Total Time (s)} & \textbf{Aborted Nodes} \\
\midrule
ExtD\_5val & 8 & \textbf{100.0\%} & 95.59\% & 75.4s & 0 \\
9Val\_Muth & 8 & 89.71\% & \textbf{98.16\%} & 321.1s & 27 \\
ExtD\_5val & 9 & \textbf{100.0\%} & 95.59\% & 235.1s & 0 \\
9Val\_Muth & 9 & 89.34\% & \textbf{98.16\%} & 344.7s & 28 \\
\bottomrule
\end{tabular}}
\end{table}

\subsection{Verification Breakdown \& Root Finding}

\textbf{Finding 1: ExtD 5-Val Overestimation and Result Fragility} \\
When computing the deepest states structurally allowed for s298 (8 time frames and 9 time frames spanning the depth of 7), the ExtD\_5val algorithm claims a stellar $100\%$ detection. It aborted exactly 0 faults, theoretically identifying valid PI assignments for everything. However, the \textit{independent fault simulation} reveals this to be inherently fragile—it concretely covers only \textbf{95.59\%} of faults, implying multiple tests returned by ExtD implicitly relied on unresolved combinations masking states in simulation. 

\textbf{Finding 2: 9Val Muth's High Veracity Models} \\
By stark contrast, Muth's 9-valued strategy inherently preserves uncertainty dynamically. It acts far more conservative when analyzing the s298 constraints over 8 frames, actively recognizing its limitations and aborting roughly $\sim$25 variables while claiming only an 89.71\% coverage. Astonishingly, when passed into concrete verification simulation, its returned test models correctly detected \textbf{98.16\%} of all faults natively! Its robustness substantially exceeds its own conservative estimates by strictly separating explicit independent structural limits over time logic. 

\textbf{Finding 3: Search Budget (Time and Scale Dependencies)} \\
Muth's algorithm requires approximately 3-to-4 times more computational budget processing deeply unrolled sequential graphs. Evaluating `frames=8` over s298 consumed roughly 321 seconds compared to ExtD's 75 seconds runtime baseline. This happens because maintaining decoupled values across the state vectors avoids rapid search contractions/abandonments.

\newpage
\section{Viva Preparation Guide (Q\&A)}

Use this section to prepare dynamically for group vivas. These encapsulate the critical concepts and findings evaluated throughout development.

\subsubsection*{Q1: Explain what exactly an "Unknown" (X) state represents in our algorithm?}
\textbf{Answer:} An 'X' state fundamentally represents lack of information. It means we cannot definitively assign a 1 or 0 yet. In sequential logic, it strictly affects flip-flop initializations. Because there's no immediate electrical pulse proving what state a D-register holds on bootup, it acts mathematically as arbitrary paths blocking combinational routes forward unless properly initialized by earlier timeframe components.

\subsubsection*{Q2: Why does the 5-Valued Algorithm suffer from overestimation and false positives in our data?}
\textbf{Answer:} Real logic relies on rigid concrete paths. ExtD_5val collapses complex signals generically—an `X` mapped structurally acts identically across both the good and faulty circuitry branches. This means the engine "guesses" paths through sequential logic assuming those overlapping `X` states won't destruct identically on combinations downstream. When the Fault Simulator aggressively binds those bounds realistically via zeros, those assumptions break resulting in lost vector coverages.

\subsubsection*{Q3: Specifically, what does 9-Valued Logic do that prevents this?}
\textbf{Answer:} Muth decoupled the model into an independent pair format \texttt{(good=0/faulty=X)} inside \texttt{logic9val.py} (\texttt{encode}/\texttt{decode} at lines 27-34). Now, if the objective determines the good route mathematically requires a 0-state, the engine validates combinations perfectly across the forward bounds even while the faulty component remains dynamically unresolvable. It natively identifies untestable topologies without guessing destructible false positives resolving vector generations inherently more resilient.

\subsubsection*{Q4: What is time frame expansion (unrolling) and why is it based on the S-Graph ($d_{\text{seq}}$)?}
\textbf{Answer:} ATPG searches are inherently combinatorial calculations—they assume signals propogate instantly. Because Sequential logic acts via cyclic clock bounds trapping bits across arrays iteratively, we clone identical blocks of combinatorial layouts referencing lines 72-139 inside \texttt{timeframe.py}. By mapping DFF-outputs from frame 't-1' directly to DFF-inputs on frame 't', we flatten time structurally. The \textit{Sequential Depth ($d_{\text{seq}}$)} found in \texttt{sgraph.py (Line 97)} proves geometrically the longest string of interacting flip-flops within circuitry indicating the bare minimum number of unrolled copies required to theoretically configure the deepest nodes mathematically possible completely. 

\subsubsection*{Q5: How exactly did you implement a Topological Sort and why was it needed?}
\textbf{Answer:} The topological parse orders generic node execution exclusively ensuring that component $B$ only ever simulates electrically once component $A$ (its parent feeder driver) structurally passes parameters recursively. We do this inside \texttt{parser.py} cleanly around \texttt{Line 174} utilizing Kahn's BFS. Standard dict tracking fails guaranteeing valid path calculations risking infinite loop recursions if backward traces evaluate before parents natively execute.

\subsubsection*{Q6: We saw Backtrack Limits scaled up from 10 to 50 in our DSE CSV. Did it help on \texttt{s298}?}
\textbf{Answer:} Scarcely. Modifying execution caps inside \texttt{atpg\_9val.\_search (Line 144)} increases computational scaling exponentially allowing the ATPG DFS algorithm additional budgets evaluating node paths dynamically. However, since \texttt{s298} features deep sequential constraints blocking true topological resolutions natively under limited time limits, additional backtracks overwhelmingly just waste execution iterations confirming dead loops eventually abandoning runs identically anyway.

\subsubsection*{Q7: Explain the D-Frontier.}
\textbf{Answer:} It is the tip of the spear during signal propagation evaluated continuously (found inside \texttt{\_d\_frontier} at \texttt{atpg\_9val L402}). It is a mathematical set tracking active gates containing an isolated discrepancy sequence (D or D\_BAR) directly applied to one or more input pins but where the combined gate-level mathematical output remains an undetermined (X) state preventing evaluation mapping to a Primary Output. Structural propagation occurs specifically by assigning non-controlling parameters driving constraints unlocking D-Frontier bounds recursively forward.

\newpage
\section{Conclusion}

This project implemented two of the foundational verification tools deployed traditionally around ASIC layouts. Beyond just theoretical metrics, analyzing an in-depth design space exposed the critical differences across sequential test coverage patterns: 

\begin{enumerate}
    \item For shallow networks like `s27`, both baseline D algorithms effectively map variables successfully ($100\%$ matched efficiency mapping). ExtD easily bridges the structural requirements within minimal steps (fast and identical output traces relative to 9-Val).
    \item For deeply layered networks like `s298`, ExtD\_5Val outputs structurally fragile test cases that simulate inconsistently in practice natively across verification scripts inside \texttt{fault\_sim.py}. 
    \item Muth's 9-valued states predictably correct this dynamic, accurately projecting rigorous verifications that resolve in significantly stronger test vectors under physical constraints (98.16\% accuracy achieved versus 95.59\%). 
    \item The trade-off penalty specifically focuses on algorithmic runtime bounds mapped internally throughout \texttt{main.py}, where 9-Val draws proportionally larger computational timeframes evaluating deeply unrolled memory chains globally.
\end{enumerate}

\appendix
\section{Appendix: Core Algorithm CLI Operations}

The codebase utilizes \texttt{main.py} dynamically mapping argument traces across modular execution patterns internally avoiding complex shell script dependencies native to group verifications:

\begin{lstlisting}[language=bash, caption=CLI Command Parameters]
# Standard Dual comparison evaluating `s27`
python main.py

# Specify individual testing arrays actively
python main.py --bench benchmarks/s298.v --mode ext_d --frames 8 --bt 50

# Output intensive design array producing design_space.csv mapped identically
python main.py --explore

# Generate Interactive HTML Graphical Circuit Components internally
python src/visualize.py benchmarks/s27.v --out results/vis
\end{lstlisting}

\section{Glossary Summary}
\begin{table}[H]
\centering
\resizebox{\textwidth}{!}{
\begin{tabular}{lp{12cm}}
\toprule
\textbf{Term} & \textbf{Definition} \\
\midrule
ATPG & Automatic Test Pattern Generation — logical combinations generating input patterns forcing exposures. \\
Stuck-At-0/1 & Persistent fault masking combinations locking wire paths independently regardless of structural parent logic bounds electrically. \\
D ($1/0$) & Multi-value signal indicating standard healthy circuits possess standard 1 voltages actively restricted to 0 bounds natively by systemic topology failure strings. \\
D-Frontier & Set of active elements blocking output combinations via undetermined evaluations natively holding active discrepancies securely inside bounds. \\
Backtracking & Recursive tracking reversal—removing assignment states entirely replacing path combinations dynamically avoiding conflict boundaries statically. \\
Time Frame Unrolling & Converting iterative memory nodes structurally mapped mathematically replacing array feedback loops actively across independent block clones dynamically linked explicitly. \\
Sequential Depth ($d_{\text{seq}}$) & Geometrical limitation identifying longest topological graph lengths mapped directly separating dependent boundary loops. \\
PPI & Pseudo-Primary Input — Earliest flip-flop signal array mathematically unbound strictly configuring global evaluation combinations passively. \\
Controlling Value & Logic parameter isolating output states uniquely locking signal boundaries exclusively (0 logic in AND / 1 logic via OR scopes). \\
\bottomrule
\end{tabular}}
\end{table}

\end{document}
